\textbf{Актуальность работы.}

Данная работа посвящена исследованию и имитационному моделированию ресурсных систем массового обслуживания. Рассматриваемая в работе модель $GI/G/N/N/R$ является частным случаем класса многосерверных систем с ограниченными ресурсами и интенсивностями. Фокус на системе с потерями позволяет явно изучить чувствительность к блокировкам по объему ресурса $R$ и количеству серверов $N$. 

В связи со сложностью аналитической оценки такой СМО, в ходе работы используется имитационное дискретно-событийное моделирование для эмпирической проверки на имитации влияния разных интенсивностей входящего потока и разных режимов распределения времени работы при фиксированном $\lambda$ на вероятность отказа, пропускную способность, среднее число заявок в системе, средний занятый ресурс, среднее время обслуживания и среднее время пребывания заявки в системе.

\textbf{Цель работы.} 

Целью работы является исследование чувствительности стационарных характеристик ресурсной системы массового обслуживания с потерями заявок $GI/G/N/N/R$ к распределению времени обслуживания и типу входящего потока при фиксированных средних значениях соответствующих случайных величин.

\textbf{Задачи работы.}

Для достижения поставленной цели в работе решаются следующие задачи: \begin{enumerate} 
    \item сформулировать математическую модель ресурсной системы массового обслуживания с потерями с ограничениями по числу обслуживающих приборов и по суммарному ресурсу; 
    \item определить набор исследуемых стационарных характеристик; 
    \item задать семейства распределений входящего потока и объёма работы заявки с одинаковыми первыми моментами, но различными коэффициентами вариации; 
    \item разработать алгоритм имитационного моделирования; 
    \item реализовать программно-вычислительный комплекс для запуска разработанного алгоритма; 
    \item провести вычислительные эксперименты и проанализировать влияние распределения времени обслуживания, типа входящего потока и их совместного действия на основные показатели эффективности системы. 
\end{enumerate}

\textbf{Объект и предмет исследования.} 
Объектом исследования являются ресурсные системы массового обслуживания с потерями, в которых поступающие заявки занимают обслуживающий прибор и случайный объём ограниченного ресурса. 

Предметом исследования является чувствительность стационарных характеристик модели $GI/G/N/N/R$ к распределению межприходных интервалов и распределению объёма работы заявки.

\textbf{Методы исследования.} 

В работе используются методы теории массового обслуживания, методы имитационного дискретно-событийного моделирования, метод Монте-Карло, методы статистической обработки результатов независимых репликаций, а также методы вычислительного эксперимента. Для оценки статистической надёжности результатов применяются выборочные средние, выборочные дисперсии, стандартные ошибки и доверительные интервалы. 

\textbf{Практическая значимость.} 

Практическая значимость работы состоит в разработке программного комплекса для имитационного анализа ресурсных систем массового обслуживания с потерями. Реализация позволяет задавать различные семейства распределений входящего потока и времени обслуживания, запускать серии независимых репликаций, агрегировать результаты и строить графики чувствительности. Полученные результаты могут быть использованы при анализе телекоммуникационных и вычислительных систем. 

\textbf{Структура работы.}

Работа состоит из введения, трех глав, заключения и приложения.
Первая глава включает обзор подходов и постановка задачи исследования» носит теоретико-аналитический характер. 
Вторая глава посвящена математическому моделированию системы, логики работы имитации и исследуемым показателям. В ней строго заданы целевые метрики (вероятность отказа, пропускная способность, средний занятый ресурс, структура состояний $\pi(k)$), описан дискретно-событийный алгоритм метода Монте-Карло и обоснован выбор семейств распределений. 
Третья глава посвящена анализу и визуализации полученных в ходе запуска имитации данных. Последовательно проверяется чувствительность к распределению времени обслуживания, затем к распределению входящего потока заявок и наконец взаимное влияние разных распределений на ключевые показатели. 

Архитектура разработанного программного комплекса, включает:
\begin{itemize}
    \item \texttt{Python} задает конфигурацию серий экспериментов (файлы \texttt{values.py}, \texttt{values\_validation.py}, \texttt{export\_values.py}) и визуализацию параметров чувствительности (\texttt{plots.py}, \texttt{plots\_reworked.py}).
    \item \texttt{Rust} модуль отвечает за запись параметров в \texttt{ExperimentConfig}, построение декартовых произведений множеств для генерации сетки сценариев (\texttt{scenario\_grid.rs}), безопасную оркестрацию параллельных потоков (\texttt{experiments.rs}) и итоговую статистическую агрегацию результатов (\texttt{stats.rs}, \texttt{output.rs}).
    \item \texttt{CUDA (C/C++)} используется  для реализация loss-системы в виде высокопараллельного потока сценариев на GPU ядрах (\texttt{sim\_kernel.cu}).
\end{itemize}

В \textbf{Заключении} обобщены итоги проведенного исследования и сформулированы основные выводы, подтверждающие вынесенные на защиту научные и инженерные гипотезы. В \textbf{Приложении А} приведены вспомогательные таблицы конфигураций и базовых параметров вычислительных экспериментов.